\section{Step 15: Compare the finalized checksum}
When both input windows are exhausted, complement the low 32 bits by XORing
with \texttt{0x00000000ffffffff}. Read the stored checksum as an unsigned
big-endian integer and compare it with the finalized accumulator. Equality
establishes that these bytes have the expected CRC; it does not establish that
the payload commands are meaningful.

Fill \texttt{record\_crc\_check}. A mismatch follows the failure path without
writing the image or context. The public \texttt{late-skip} input has a freshly
computed, correct CRC despite an invalid command. It will pass this region and
must still fail during command validation.

\paragraph{Observe} Prepare \texttt{s15-1} and stop before
\texttt{record\_crc\_checked}. At \texttt{record\_crc\_checked}, both \texttt{\$t2} and
\texttt{\$t8} are \texttt{0x00000000ba0fd49e} for the worked record.
The \texttt{bad-crc} recipe returns one before reaching this checkpoint.

\Needspace{7\baselineskip}
\paragraph{Check} Does a matching CRC prove command bounds are valid?
\begin{itemize}
  \item[A.] Yes, because CRCs check the format.
  \item[B.] Yes, if the record is short.
  \item[C.] No; command bounds need separate checks.
\end{itemize}
